
\begin{theorem}[Conditional UV Bernstein-cone program]
\label{thm:uv-bernstein-cone-completion}
Assume the projected RG flow preserves the admissible Bernstein cone and admits a nontrivial UV
fixed point satisfying the pole and positivity audit below. The effective average action is
sought in the form
\[
\Gamma_k
=
\int d^4x\sqrt g\,
\Bigl[
\frac{M_k^2}{2}R
+C\,\mathcal F_{C,k}(\Box/\chigeo^2)\,C
+R\,\mathcal F_{R,k}(\Box/\chigeo^2)\,R
+\frac12 Z_{\chi,k}(\nabla\chigeo)^2
+U_k(\chigeo,\Phi)
+\cdots
\Bigr],
\]
where the sector form factors are constrained by complete Bernstein functions
\[
\Phi_{s,k}(z)
:=
a_{s,k}+b_{s,k}z+\int_0^\infty \frac{z}{u+z}\rho_{s,k}(u)\,du,
\qquad
\rho_{s,k}(u)\ge 0.
\]
The admissible RG flow is required to preserve both the retained sector and this
cone:
\[
\partial_t\Gamma_k
:=
\frac12\mathrm{STr}_{\Padm}
\Bigl[
\bigl(\Gamma_k^{(2)}+R_k\bigr)^{-1}\partial_t R_k
\Bigr],
\]
with regulator $R_k$ commuting with $\Padm$, $Q_\Theta$, and $\inv$ and living in
the same Bernstein class. Sectorwise the inverse propagators are then required to
take the form
\[
\Gamma^{(2)}_{2,k}(p^2)
=
Z_{N,k}\,p^2\,\Phi_{2,k}(p^2/\chigeo^2),
\]
\[
\Gamma^{(2)}_{0,k}(p^2)
=
Z_{0,k}\,(p^2+M_{s,k}^2)\,\Phi_{0,k}(p^2/\chigeo^2).
\]
The conditional UV program closes only if three conditions hold simultaneously:
\begin{enumerate}[label=(\roman*), leftmargin=1.5em, topsep=3pt, itemsep=2pt]
    \item the admissible Bernstein cone is invariant under the projected RG flow,
    \item the flow admits a nontrivial UV fixed point on that cone,
    \item and the pole/positivity audit closes on $\operatorname{Ran}(\Padm)$:
    exactly one massless spin-$2$ pole, exactly one healthy scalaron pole with
    $M_{s,k}^2>0$ and positive residue, no extra ghost/tachyon branches, and
    positive spectral density compatible with admissible reflection positivity.
\end{enumerate}
Under these assumptions, the projected RG flow preserves the admissible complete-Bernstein cone,
admits a nontrivial UV fixed point on the bounded moment body, and closes the pole/positivity
audit on $\operatorname{Ran}(\Padm)$ with exactly one massless spin-$2$ pole, exactly one
healthy scalaron pole, and no extra ghost or tachyon branches. The local Einstein plus $R^2$
branch is therefore interpreted as the infrared projection of an admissible UV Bernstein-cone
program.
\end{theorem}

\begin{theorem}[Conditional regulator equivalence inside the admissible Bernstein cone]
\label{thm:regulator-equivalence-bernstein-cone}
Assume the hypotheses of \Cref{thm:uv-bernstein-cone-completion}. Let $R_k$ and
$\widetilde R_k$ be two admissible projected regulators in the same complete Bernstein class,
both commuting with $\Padm$, $Q_\Theta$, and $\inv$, and both normalized to the same infrared
data. Then their induced admissible RG flows determine the same:
\begin{enumerate}[label=(\roman*), leftmargin=1.5em, topsep=3pt, itemsep=2pt]
    \item fixed-point pole data in the spin-$2$ and spin-$0$ sectors,
    \item number of relevant UV directions,
    \item infrared Einstein--$R^2$ closure class on the admissible branch.
\end{enumerate}
Hence the conditional UV program of \Cref{thm:uv-bernstein-cone-completion} is regulator
independent inside the admissible Bernstein-cone class.
\end{theorem}

The E$_8$ scale ladder is not used as a primitive input in the main text. Only its short downstream grammar is retained in the dedicated main-text section, while the fuller arithmetic continuations are deferred to Appendix~\ref{app:constants}.

\section{Closure architecture B: Yukawas, masses, flavor, confinement, and strong CP}

The second closure-architecture block is the transport compiler. Its role is to turn the carrier packet and the family/occupancy closure into actual masses, mixings, hadronic admissibility, and strong-CP control. The logic is again staged: first the discrete phase operator and the transport kernel, then path-length masses and mixing data, then confinement and strong CP as transport-side admissibility statements.

\subsection{Discrete phase operator and transport grammar}

Let
\[
\omega=e^{2\pi i/3},
\qquad
U_6=\mathrm{diag}(1,\omega,\omega^2,-1,-\omega,-\omega^2).
\]
This makes the phase grammar concrete rather than rhetorical. The six admissible phases are shown as a finite hexagonal transport geometry in \Cref{fig:z6-circle}: the cusp values live on the same real axis as the positive-real phase, the admissible gap is visible directly, and the holonomy phase can be read as a triangle on the same discrete transport base.

\begin{remark}[Hexagonal spectral geometry]
The operator $U_6$ is a unitary circulant with eigenphases $e^{in\pi/3}$ for $n=0,\dots,5$ and is unitarily equivalent to the one-step translation operator on a six-site ring. The transport kernel $D_y=y\,\mathbf 1-\delta_{\mathrm{ph}}U_6$ is therefore the resolvent denominator of a finite quantum walk, equivalently a tight-binding transfer problem, on the discrete $C_6$ cycle. The regularized inverse introduced below plays the role of the corresponding Green kernel. CKM and PMNS mixings, mass hierarchies, and the Möbius quotients of \Cref{eq:mobius-quotient} are thus not loose analytic fits but spectral-theory statements on a regular hexagon. The near-pole hierarchy at $y\approx\delta_{\mathrm{ph}}$ becomes the resonance of a hexagonal resolvent rather than tuned parameter choice.
\end{remark}

\begin{figure}[t]
\centering
\scriptsize
\begin{tikzpicture}[scale=1.1,>=Latex]
\coordinate (O) at (0,0);
\foreach \name/\ang in {p0/0,p1/60,p2/120,p3/180,p4/240,p5/300} {
    \coordinate (\name) at ({2*cos(\ang)},{2*sin(\ang)});
}
\draw[thick, black!70] (p0) -- (p1) -- (p2) -- (p3) -- (p4) -- (p5) -- cycle;
\draw[gray!25, dashed] (O) -- (p0);
\draw[gray!25, dashed] (O) -- (p1);
\draw[gray!25, dashed] (O) -- (p2);
\draw[gray!25, dashed] (O) -- (p3);
\draw[gray!25, dashed] (O) -- (p4);
\draw[gray!25, dashed] (O) -- (p5);
\foreach \P in {p0,p1,p2,p3,p4,p5} {
    \filldraw[blue!60!black] (\P) circle (2pt);
}
\node[anchor=west, font=\scriptsize] at (2.18,0.12) {$1$};
\node[anchor=south west, font=\scriptsize] at (1.1,1.85) {$e^{i\pi/3}$};
\node[anchor=south east, font=\scriptsize] at (-1.1,1.85) {$e^{2i\pi/3}$};
\node[anchor=east, font=\scriptsize] at (-2.18,0.12) {$-1$};
\node[anchor=north east, font=\scriptsize] at (-1.1,-1.85) {$e^{-2i\pi/3}$};
\node[anchor=north west, font=\scriptsize] at (1.1,-1.85) {$e^{-i\pi/3}$};

\draw[->, thick, red!65!black] (1.15,1.05) arc[start angle=42,end angle=-42,radius=1.2];
\node[red!65!black, font=\scriptsize, anchor=west] at (2.55,0.25) {CP: $\theta\mapsto-\theta$};

\filldraw[orange!80!black] (0.67,0) circle (2pt);
\filldraw[orange!80!black] (1.33,0) circle (2pt);
\filldraw[green!60!black] (1.1998,0) circle (2.4pt);
\node[below, font=\scriptsize] at (0.67,-0.08) {$\frac13$};
\node[below, font=\scriptsize] at (1.33,-0.08) {$\frac23$};
\node[above, font=\scriptsize] at (1.1998,0.08) {$\delta_{\mathrm{ph}}$};
\draw[<->, thick, black!65] (1.1998,-0.32) -- (1.33,-0.32);
\node[below, font=\tiny] at (1.265,-0.34) {$\epsilon_f$};
\node[align=left, anchor=west, font=\scriptsize, text=black!70] at (2.55,-0.7)
    {near pole at\\$y\!\approx\!\delta_{\mathrm{ph}}$\\largest hierarchy gain};

\begin{scope}[shift={(4.6,-0.8)}]
    \coordinate (A) at (0,0.9);
    \coordinate (B) at (-0.85,-0.5);
    \coordinate (C) at (0.85,-0.5);
    \draw[thick, orange!75!black] (A) -- (B) -- (C) -- cycle;
    \filldraw[orange!75!black] (A) circle (1.5pt) (B) circle (1.5pt) (C) circle (1.5pt);
    \node[above, font=\scriptsize] at (A) {$1$};
    \node[left, font=\scriptsize] at (B) {$2$};
    \node[right, font=\scriptsize] at (C) {$3$};
    \node[align=center, font=\scriptsize, text=black!70] at (0,-1.1) {$\delta_{\mathrm{CKM}}=\arg\tr\Pi_\triangle$\\holonomy triangle};
\end{scope}
\end{tikzpicture}
\caption{Finite transport geometry on the hexagon $C_6$. The six $U_6$ phases form a discrete transport base, the admissible cusps $y\in\{1,\frac13,\frac23\}$ sit on the positive-real axis together with $\delta_{\mathrm{ph}}$, the smallest gap is $\epsilon_f=|\frac23-\delta_{\mathrm{ph}}|$, and the same transport grammar also carries the holonomy triangle phase.}
\label{fig:z6-circle}
\end{figure}

\subsection{Regularized Yukawa kernels and path-length masses}

\[
P_{\mathrm{str}}:=P_{\Seam,+}P_{\mathrm{sing}}P_\Theta,
\qquad
P_{\Seam,+}:=\frac{\mathbf 1+\inv}{2},
\qquad
P_\Theta:=\frac12(\mathbf 1+Q_\Theta),
\qquad
P_{\mathrm{sing}}:=\frac13(\mathbf 1+Z_c+Z_c^2).
\]

\begin{definition}[Structural transport kernel]
For each fermion sector define the structural positive transport kernel by
\[
K_f^{\mathrm{str}}:=\bigl(P_{\mathrm{str}} D_f^\dagger D_f P_{\mathrm{str}}\bigr)^{-1/2},
\]
whenever the inverse exists on $\mathrm{Ran}(P_{\mathrm{str}})$. The final admissible kernel is
the restriction of this operator to the gap-stable branch selected below. The transport pole
$\delta_{\mathrm{ph}}$ is the unique stationary point of the exact determinant functional
\[
\mathcal F_{\mathrm{tr}}(\delta)
:=
\sum_{y\in\{1,1/3,2/3\}}
\Bigl[-\log\det\!\bigl(D_{y,\delta}^\dagger D_{y,\delta}+\varepsilon^2\bigr)\Bigr]
+
\frac{36}{(\phiz)^2}\left(\delta-\frac35\right)^2,
\]
with
\[
D_{y,\delta}:=y\,\mathbf 1-\delta U_6.
\]
On the canonical admissible branch one has
\begin{equation}
\label{eq:transport-bridge-offset}
\partial_\delta \mathcal F_{\mathrm{tr}}(\delta_{\mathrm{ph}})=0,
\qquad
\delta_{\mathrm{ph}}\in\left(\frac13,\frac23\right),
\qquad
\delta_{\mathrm{ph}}\approx \num{0.599882}.
\end{equation}
\end{definition}

\begin{theorem}[Canonical admissible transport branch]
\label{thm:exact-transport-closure}
Assume $\mathcal V_{\mathrm{vac}}$ is coercive on the retained admissible branch and strictly
convex in $\delta$ and in the chosen holonomy coordinates near the selected minimizer.
Define the global admissible vacuum functional
\[
\mathcal V_{\mathrm{vac}}[g,\chigeo,\Phi,\nabla_F,\delta,\{L_f\}]
:=
\Gamma_{\mathrm{geom}}[g,\chigeo,\Phi]
+\mathcal F_{\mathrm{tr}}(\delta)
-\sum_f \log \mathcal O_f(L_f,\nabla_F)
+\mu\sum_{f,j}L_{f,j}
+\Lambda_\Seam^{\mathrm{seam}},
\]
with exact overlaps
\[
\mathcal O_f(L,\nabla_F):=
\langle\Omega|
T_f(L,\nabla_F)^\dagger
K_f^{\mathrm{str}}
T_f(L,\nabla_F)
|\Omega\rangle.
\]
Then there exists a unique local admissible minimizer
\[
(\delta_{\mathrm{ph}},\nabla_F^\star,\{L_f^{\mathrm{diag}}\}).
\]
In the present manuscript the packet triples
\[
L_u^{\mathrm{diag}}=(7,3,0),\quad
L_d^{\mathrm{diag}}=(7,5,2),\quad
L_e^{\mathrm{diag}}=(8,5,3)
\]
are the selected values on that canonical admissible branch, and the Yukawa matrices are
\[
(Y_f)_{ij}
=
\sum_{C:i\to j}\lambda_Y^{\ell(C)}
\langle e_j,K_f^{\mathrm{str}}\,\mathrm{Hol}_{C}(\nabla_F^\star)e_i\rangle.
\]
\end{theorem}

Here $\delta_{\mathrm{ph}}\approx \num{0.599882}$ lies slightly below $3/5$, so $\epsilon_f=|\frac23-\delta_{\mathrm{ph}}|\approx \num{0.06678}$. The regularized transport operator used for the finite hexagon analysis is
\begin{equation}
\label{eq:yukawa-kernel}
\begin{aligned}
D_y&=y\,\mathbf 1-\delta_{\mathrm{ph}}U_6,\\
\mathcal Y_y^{(\varepsilon)}&=(D_y^\dagger D_y+\varepsilon^2)^{-1},
\end{aligned}
\end{equation}
and the singular limit $\varepsilon\to0$ is taken only after admissibility has been restricted to the gap-stable sector defined below.

\begin{proposition}[Finite hexagonal resolvent]
\label{prop:finite-hexagonal-resolvent}
Whenever
\[
y^6\neq \delta_{\mathrm{ph}}^6,
\]
the transport denominator has the exact inverse
\begin{equation}
\label{eq:finite-hexagonal-resolvent}
D_y^{-1}
=
\frac{1}{y^6-\delta_{\mathrm{ph}}^6}
\sum_{m=0}^{5}y^{5-m}\delta_{\mathrm{ph}}^{\,m}U_6^m.
\end{equation}
Thus the $C_6$ resolvent is finite and algebraically closed rather than an infinite formal Green-series expansion.
\end{proposition}

\begin{proof}
Using $U_6^6=\mathbf 1$,
\[
\bigl(y\,\mathbf 1-\delta_{\mathrm{ph}}U_6\bigr)
\sum_{m=0}^{5}y^{5-m}\delta_{\mathrm{ph}}^{\,m}U_6^m
=
y^6\mathbf 1-\delta_{\mathrm{ph}}^6U_6^6
=
\bigl(y^6-\delta_{\mathrm{ph}}^6\bigr)\mathbf 1.
\]
Dividing by $y^6-\delta_{\mathrm{ph}}^6$ gives \Cref{eq:finite-hexagonal-resolvent}.
\end{proof}

Equivalently, one may read \Cref{eq:yukawa-kernel} in the Fourier basis of the six-site ring: $U_6$ is the diagonal form of the step operator, $D_y$ is the shifted walk denominator, and $\mathcal Y_y^{(\varepsilon)}$ is the corresponding regularized Green function. In this reading the admissible path lengths $L_{f,j}$ are graph-distance data on $C_6$, while the transport matrix elements $\Lambda_{f,j}$ encode bounded local amplitudes rather than unrestricted fit parameters.

Because $U_6$ has eigenphases $e^{in\pi/3}$, the regularized transport spectrum is
\begin{equation}
\label{eq:transport-spectrum}
\lambda_{y,n}^{-1}
=
y^2+\delta_{\mathrm{ph}}^2-2y\delta_{\mathrm{ph}}\cos\!\left(\frac{n\pi}{3}\right)+\varepsilon^2,
\qquad
n=0,\dots,5,
\end{equation}
for the eigenvalues $\lambda_{y,n}$ of $\mathcal Y_y^{(\varepsilon)}$. In particular,
\begin{lemma}[Finite admissible hexagon gap]
\label{lem:finite-admissible-hexagon-gap}
On the admissible cusp set
\[
y\in\left\{1,\frac13,\frac23\right\},
\]
the smallest singular value of the unregulated hexagon denominator is attained at $y=\frac23$ and $n=0$. Equivalently,
\[
\epsilon_f
:=
\min_{\substack{y\in\{1,\frac13,\frac23\}\\ n=0,\dots,5}}
\left|y-\delta_{\mathrm{ph}}e^{in\pi/3}\right|
=
\left|\frac23-\delta_{\mathrm{ph}}\right|.
\]
Hence the finite admissible hexagon has a strictly positive spectral gap $\epsilon_f$.
\end{lemma}

\begin{proof}
For $\varepsilon=0$, the eigenvalues in \Cref{eq:transport-spectrum} reduce to $|y-\delta_{\mathrm{ph}}e^{in\pi/3}|^{-2}$. On the retained branch one has $\delta_{\mathrm{ph}}\approx \num{0.599882}$, so among the admissible cusps the closest point to the positive real phase is $y=2/3$ with $n=0$. This gives the stated minimum.
\end{proof}

\begin{definition}[Gap-stable admissibility]
\label{def:gap-stable-admissibility}
Let
\[
D_0:=D_{2/3}.
\]
We say that the admissible transport sector is gap-stable if the physical transport operator has the form
\[
D_f=D_0+V_{\mathrm{adm}},
\qquad
[P_{\mathrm{str}},V_{\mathrm{adm}}]=0,
\qquad
\|V_{\mathrm{adm}}\|<\frac{\epsilon_f}{2}.
\]
\end{definition}

\begin{proposition}[Structural gap lift]
\label{prop:gap-lift-full-operator}
If admissibility is gap-stable in the sense of \Cref{def:gap-stable-admissibility}, then for every
\[
\psi\in\mathrm{Ran}(P_{\mathrm{str}})
\]
one has
\[
\|D_f\psi\|
\ge
\frac{\epsilon_f}{2}\|\psi\|.
\]
Consequently
\begin{equation}
\label{eq:gap-lift-full-operator}
P_{\mathrm{str}} D_f^\dagger D_f P_{\mathrm{str}}
\ge
\frac{\epsilon_f^2}{4}\,P_{\mathrm{str}}
>
0.
\end{equation}
In particular, the final admissible restriction to $\mathrm{Ran}(\Padm)$ is strictly positive.
\end{proposition}

\begin{proof}
Since $\psi\in\mathrm{Ran}(P_{\mathrm{str}})$ and $[P_{\mathrm{str}},V_{\mathrm{adm}}]=0$,
\[
\|D_f\psi\|
\ge
\|D_0\psi\|-\|V_{\mathrm{adm}}\psi\|
\ge
\bigl(\epsilon_f-\|V_{\mathrm{adm}}\|\bigr)\|\psi\|
\ge
\frac{\epsilon_f}{2}\|\psi\|.
\]
Squaring gives \Cref{eq:gap-lift-full-operator}.
\end{proof}

\begin{equation}
\label{eq:mobius-quotient}
\frac{\lambda_{y,0}}{\lambda_{y,3}}
=
\frac{(y+\delta_{\mathrm{ph}})^2+\varepsilon^2}{(y-\delta_{\mathrm{ph}})^2+\varepsilon^2}.
\end{equation}
\[
\lim_{\varepsilon\to 0}\frac{\lambda_{y,0}}{\lambda_{y,3}}
=
\left(\frac{y+\delta_{\mathrm{ph}}}{y-\delta_{\mathrm{ph}}}\right)^2
\]
\begin{remark}[Near-pole hierarchy mechanism]
On the retained canonical transport saddle from \Cref{eq:transport-bridge-offset},
\[
\delta_{\mathrm{ph}}\approx \num{0.599882}.
\]
Hence the Möbius quotient contains a built-in near-pole whenever $y$ approaches $\delta_{\mathrm{ph}}$, making large hierarchy ratios structurally plausible. In the closed transport grammar this is the hierarchy mechanism, while the bounded holonomy factors keep it from degenerating into an arbitrary tuning surface.
\end{remark}
